\section{Motivation}
\label{sec:motivation}

The rapid growth of digital content platforms has fundamentally changed how users discover and consume information.
In the domain of books alone, major platforms now catalogue tens of millions of titles, making the curation of relevant recommendations a non-trivial problem at scale.
A user browsing a catalogue of three million books faces a search space that manual exploration cannot realistically traverse; a recommendation system must therefore act as an intelligent intermediary that surfaces relevant items before the user expresses a precise query.

Traditional recommendation approaches tend to be unimodal.
Collaborative filtering methods exploit the co-occurrence of user–item interactions but ignore the semantic content of the items themselves.
Content-based methods encode item descriptions into dense vectors and retrieve semantically similar content, but are insensitive to the community-level purchasing patterns that often drive discovery.
Neither class of method captures \textit{sequential intent} — the fact that a user's most recent interactions often signal a current reading interest that differs from their long-term profile.
This gap motivates the development of a \textit{dual-mode} architecture that combines both behavioural signals and sequential intent modelling within a unified recommendation pipeline.

A further practical challenge is language.
In Vietnam and other multilingual markets, a substantial portion of users formulate queries in their native language, yet most open-source encoders are trained primarily on English text and perform poorly on Vietnamese input.
This discrepancy creates a de facto exclusion of non-English speakers from the quality of search results that English speakers receive.
Addressing multilingual retrieval is therefore not merely a technical refinement but a question of equitable access.

Finally, real-world deployment demands that accuracy improvements not come at the cost of user experience.
A recommendation API that takes more than one second to respond is unlikely to be adopted in an interactive application.
Latency optimisation — through model quantisation, approximate nearest-neighbour indexing, and response caching — must therefore be treated as a first-class design constraint rather than an afterthought.

This capstone project addresses these challenges by designing, implementing, and evaluating a multimodal recommendation system that combines Cleora co-purchase graph embeddings with a DIF-SASRec sequential transformer, supported by a multilingual BGE-M3 encoder and a latency-optimised search pipeline.
The system is built on a modular data platform whose ingestion, storage, and visualisation components are additionally demonstrated on an air quality monitoring use case, establishing the generality of the underlying infrastructure.

